\scanpage{28}
\mspage{15}

$Y_{J,K}\to\deleted{\unsure{X}}Y_{J,L}$. Nous utiliserons ce diagramme dans les deux cas où $(J\subset K\subset L)$ est l’un de deux triples

\[
 \left\{\begin{aligned}
 I'_0&\subset I_0\subset I\\
 I'_0&\subset I'\subset I
 \end{aligned}\right.
\]

\inserted{issus} du carré cocartésien (\deleted{\unsure{de la}} cat) (\deleted{\unsure{ou}} carré MV)

\diagram{\begin{tikzcd}[ampersand replacement=\&]
 \&I'_0\arrow[r,hook,no head]\arrow[d,hook]\&I_0\arrow[d,hook]\\
 (3)\&I'\arrow[r,hook]\&I\, .
\end{tikzcd}}

Il faut insérer les deux diagrammes correspondants \unsure{du|au} type (2) dans un même diagramme \unsure{!|?} On notera à ce propos que les carrés \inserted{\unsure{de|dans} $\unsure{\mathcal M|\mathscr M}$,} déduits de (3) en lui appliquant \deleted{$\unsure{J}$} $J\mathbin{\unsure{\mapstochar\relbar|\mapsto}}X_J$ et $J\mathbin{\mapstochar\relbar}Y_J$

\diagram{$\left.\vcenter{\hbox{\begin{tikzcd}[ampersand replacement=\&,column sep=small,row sep=small]
 \&X_{I'_0}\arrow[rr]\arrow[dr]\arrow[dd,dashed,"u_{I'_0}"]\&\&X_{I_0}\arrow[dr]\arrow[dd,dashed,no head,"u_{I_0}"]\&\\
 (4_x)\&\&X_{I'}\arrow[rr]\arrow[dd,dashed,"u_{I'}"']\&\&X_I\arrow[dd,dashed,"u_I"']\\
 \&Y_{I'_0}\arrow[rr]\arrow[dr]\&\&Y_{I_0}\arrow[dr]\&\\
 (4_y)\&\&Y_{I'}\arrow[rr,no head]\&\&Y_I
\end{tikzcd}}}\right\}\ (4)$}

sont cocartésiens dans $\mathcal M$. (Transitivité des $\underrightarrow{\lim}\ldots$).
% archive sheet number: 27

